\chapter{Conclusiones}
\label{cap:conclusiones}

Esta tesis partió de una dificultad concreta: un objetivo de software puede ser
demasiado amplio para un único agente, pero dividirlo sin criterio crea más
problemas de los que resuelve. ManyHands aborda esa tensión mediante una
arquitectura en la que la semántica generativa y el control determinista cumplen
funciones distintas.

\section{Respuesta a las preguntas de investigación}

Respecto de la primera pregunta, el trabajo muestra que es viable coordinar
agentes mediante un grafo jerárquico, contratos versionados, ejecución aislada e
integración basada en evidencia. La viabilidad no se apoya sólo en el diseño:
el run experimental completó planificación, ejecución, recuperación,
integraciones intermedias, integración raíz y entrega sobre un commit exacto.

Respecto de la segunda pregunta, se formuló una política de granularidad
determinista y explicable. La decisión ya no depende de un puntaje difícil de
interpretar, sino de tres razones observables: que la tarea no quepa, que la
división habilite paralelismo o que permita verificar partes por separado. El
algoritmo de colapso conserva obligaciones y elimina fronteras sin utilidad.

\section{Aportes principales}

Los aportes del trabajo pueden resumirse en seis elementos:

\begin{enumerate}
  \item un \textbf{grafo híbrido} que separa jerarquía, flujo de artefactos,
  contratos de interfaz y conflictos de recursos;
  \item una \textbf{política de granularidad categórica} acompañada por un
  algoritmo de colapso auditable;
  \item un \textbf{catálogo de contratos} que preserva objetivo, alcance,
  artefactos y criterios durante la descomposición;
  \item un subsistema de \textbf{ejecución aislada y supervisada} sobre Git
  worktrees, sandbox y custodia del árbol de procesos;
  \item una \textbf{integración ascendente} que trata a los compuestos como
  intentos de primer orden y vincula criterios con evidencia;
  \item una \textbf{historia durable del run} y una entrega transaccional que
  relacionan el resultado publicado con su candidato físico.
\end{enumerate}

\section{Lecciones de ingeniería}

Tres ideas atravesaron el desarrollo. Primero, una declaración del modelo nunca
debe sustituir la inspección del artefacto. Segundo, los contratos reducen la
ambigüedad sólo si están conectados con mecanismos físicos de ejecución y
validación. Tercero, un sistema autónomo útil no es aquel que nunca falla, sino
el que puede localizar una falla, conservar el trabajo válido y continuar sin
perder trazabilidad.

El experimento final reunió estas ideas en un único recorrido. El resultado fue
un grafo de tres niveles, con trabajo concurrente, integraciones por área,
reparaciones localizadas, candidato verificado y entrega reproducible. Más que
una demostración visual, el valor del caso reside en que cada etapa puede
explicarse mediante eventos, commits y recibos.

\section{Trabajo futuro}

Las extensiones prioritarias son:

\begin{enumerate}
  \item incorporar oráculos de producto externos como obligaciones canónicas
  previas a la entrega;
  \item mejorar los checkpoints de intentos para reutilizar candidatos parciales
  durante una reparación;
  \item calificar el sandbox y la supervisión en más sistemas operativos;
  \item estudiar la política sobre repositorios y lenguajes diversos;
  \item evaluar usabilidad y comprensión del command center con usuarios;
  \item combinar granularidad con enrutamiento dinámico de modelos según costo y
  capacidad.
\end{enumerate}

ManyHands no elimina la incertidumbre de los modelos de lenguaje. La rodea con
una estructura de ingeniería que vuelve sus acciones observables, recuperables
y verificables. Ese es el resultado central de esta tesis y la base para futuras
plataformas de desarrollo autónomo responsables.
